\section{早明的冲突、灾害与司法证据}

战争、灾荒、疾病、赈恤、案狱和边地暴力并不是彼此独立的材料领域。军报和判牍往往由胜者、官府或审讯机关形成；灾情、医疗与慈善记录也常以请赈、捐输、工程或死亡数字的形式出现。它们能显示国家和地方组织如何命名风险与行动，却不能把受害者、病者、被征发者或被定罪者的经验全部还原。

\subsection{战事与军政暴力}

\eventanchor{ML-1385-001}{明·1385（洪武十八年）}
\paragraph{1385（洪武十八年）：吴冕叛乱：锦叛乱被击败} 此一账本锚点将冲突、风险或法律处置置于具体时段。它所能可靠支持的判断是编年候选的年序可由官修与后出材料交叉；具体月日、参与者、战果或数字必须回查所列材料，不能将单一叙事当作定论。材料链由，因而同时保存了命令、报送、传记、地方记忆或跨政权叙事的不同立场。问题形成的条件是前期边防、地方武装或跨区域资源调动的压力推动了该行动；随后产生的可见影响是它改变了后续调兵、补给或地方秩序的条件；战果和伤亡须分开核对。对于医疗救助、慈善网络、普通人的伤亡与案狱处境，材料往往尤为稀薄。译写候选以编年材料定位；实录记载反映朝廷选择，崇祯期须另以奏疏、地方及敌对政权材料交叉。

\eventanchor{ML-1385-002}{明·1385（洪武十八年）}
\paragraph{1385（洪武十八年）：科举考试恢复} 此一账本锚点将冲突、风险或法律处置置于具体时段。它所能可靠支持的判断是编年候选的年序可由官修与后出材料交叉；具体月日、参与者、战果或数字必须回查所列材料，不能将单一叙事当作定论。材料链由，因而同时保存了命令、报送、传记、地方记忆或跨政权叙事的不同立场。问题形成的条件是制度、教育或知识传播的需要推动了相关行动或文本形成；随后产生的可见影响是它提供制度和传播史的锚点，不能据此推断社会整体接受。对于医疗救助、慈善网络、普通人的伤亡与案狱处境，材料往往尤为稀薄。译写候选以编年材料定位；实录记载反映朝廷选择，崇祯期须另以奏疏、地方及敌对政权材料交叉。

\eventanchor{ML-1385-003}{明·1385（洪武十八年）}
\paragraph{1385（洪武十八年）：郭焕因贪污粮食700万担被执行死刑} 此一账本锚点将冲突、风险或法律处置置于具体时段。它所能可靠支持的判断是编年候选的年序可由官修与后出材料交叉；具体月日、参与者、战果或数字必须回查所列材料，不能将单一叙事当作定论。材料链由，因而同时保存了命令、报送、传记、地方记忆或跨政权叙事的不同立场。问题形成的条件是继承、官僚协调或皇权运作的压力使问题进入朝廷程序；随后产生的可见影响是它影响后续人事与政策议程；动机和派系标签不能由单一叙事裁定。对于医疗救助、慈善网络、普通人的伤亡与案狱处境，材料往往尤为稀薄。译写候选以编年材料定位；实录记载反映朝廷选择，崇祯期须另以奏疏、地方及敌对政权材料交叉。

\eventanchor{ML-1386-001}{明·1386（洪武十九年）一月}
\paragraph{1386（洪武十九年）一月：明蒙毛战争：蒙毛起义军的司论法} 此一账本锚点将冲突、风险或法律处置置于具体时段。它所能可靠支持的判断是编年候选的年序可由官修与后出材料交叉；具体月日、参与者、战果或数字必须回查所列材料，不能将单一叙事当作定论。材料链由，因而同时保存了命令、报送、传记、地方记忆或跨政权叙事的不同立场。问题形成的条件是前期边防、地方武装或跨区域资源调动的压力推动了该行动；随后产生的可见影响是它改变了后续调兵、补给或地方秩序的条件；战果和伤亡须分开核对。对于医疗救助、慈善网络、普通人的伤亡与案狱处境，材料往往尤为稀薄。译写候选以编年材料定位；实录记载反映朝廷选择，崇祯期须另以奏疏、地方及敌对政权材料交叉。

\eventanchor{ML-1386-002}{明·1386（洪武十九年）}
\paragraph{1386（洪武十九年）：朝廷围绕卫所军户的编籍、屯田与调发整理本年军政文书，作为明初军事接收的可核查节点} 此一账本锚点将冲突、风险或法律处置置于具体时段。它所能可靠支持的判断是来源导向的年度桥接条目：本年材料可定位相关政务线索；具体月日、发文机关、地域和执行结果仍须逐项回查，不能以制度文本或奏报替代实际效果。材料链由，因而同时保存了命令、报送、传记、地方记忆或跨政权叙事的不同立场。问题形成的条件是前期边防、地方武装或跨区域资源调动的压力推动了该行动；随后产生的可见影响是它改变了后续调兵、补给或地方秩序的条件；战果和伤亡须分开核对。对于医疗救助、慈善网络、普通人的伤亡与案狱处境，材料往往尤为稀薄。桥接条目只标示可回查的年度问题与材料链；实录有中央选择性，崇祯期尤其需要奏疏、地方、朝鲜及满文材料交叉。

\eventanchor{ML-1387-001}{明·1387（洪武二十年）十月}
\paragraph{1387（洪武二十年）十月：明朝对乌里安凯的战役：那加楚向明朝军队投降} 此一账本锚点将冲突、风险或法律处置置于具体时段。它所能可靠支持的判断是编年候选的年序可由官修与后出材料交叉；具体月日、参与者、战果或数字必须回查所列材料，不能将单一叙事当作定论。材料链由，因而同时保存了命令、报送、传记、地方记忆或跨政权叙事的不同立场。问题形成的条件是前期边防、地方武装或跨区域资源调动的压力推动了该行动；随后产生的可见影响是它改变了后续调兵、补给或地方秩序的条件；战果和伤亡须分开核对。对于医疗救助、慈善网络、普通人的伤亡与案狱处境，材料往往尤为稀薄。译写候选以编年材料定位；实录记载反映朝廷选择，崇祯期须另以奏疏、地方及敌对政权材料交叉。

\subsection{灾害、医疗与赈恤}

\eventanchor{ML-1388-001}{明·1388（洪武二十一年）五月}
\paragraph{1388（洪武二十一年）五月：贝尔湖之战：明军击败乌斯哈尔汗托古斯帖木儿} 此一账本锚点将冲突、风险或法律处置置于具体时段。它所能可靠支持的判断是编年候选的年序可由官修与后出材料交叉；具体月日、参与者、战果或数字必须回查所列材料，不能将单一叙事当作定论。材料链由，因而同时保存了命令、报送、传记、地方记忆或跨政权叙事的不同立场。问题形成的条件是前期边防、地方武装或跨区域资源调动的压力推动了该行动；随后产生的可见影响是它改变了后续调兵、补给或地方秩序的条件；战果和伤亡须分开核对。对于医疗救助、慈善网络、普通人的伤亡与案狱处境，材料往往尤为稀薄。译写候选以编年材料定位；实录记载反映朝廷选择，崇祯期须另以奏疏、地方及敌对政权材料交叉。

\eventanchor{ML-1388-002}{明·1388（洪武二十一年）}
\paragraph{1388（洪武二十一年）：明蒙毛战争：蒙毛被明军炮兵齐射击败} 此一账本锚点将冲突、风险或法律处置置于具体时段。它所能可靠支持的判断是编年候选的年序可由官修与后出材料交叉；具体月日、参与者、战果或数字必须回查所列材料，不能将单一叙事当作定论。材料链由，因而同时保存了命令、报送、传记、地方记忆或跨政权叙事的不同立场。问题形成的条件是自然风险与地方报告促成朝廷或地方的应对记录；随后产生的可见影响是赈济、迁徙和损失的实际范围仍须以受灾地区材料分别核验。对于医疗救助、慈善网络、普通人的伤亡与案狱处境，材料往往尤为稀薄。译写候选以编年材料定位；实录记载反映朝廷选择，崇祯期须另以奏疏、地方及敌对政权材料交叉。

\eventanchor{ML-1389-001}{明·1389（洪武二十二年）一月}
\paragraph{1389（洪武二十二年）一月：明军在岳州击败彝族叛军} 此一账本锚点将冲突、风险或法律处置置于具体时段。它所能可靠支持的判断是编年候选的年序可由官修与后出材料交叉；具体月日、参与者、战果或数字必须回查所列材料，不能将单一叙事当作定论。材料链由，因而同时保存了命令、报送、传记、地方记忆或跨政权叙事的不同立场。问题形成的条件是前期边防、地方武装或跨区域资源调动的压力推动了该行动；随后产生的可见影响是它改变了后续调兵、补给或地方秩序的条件；战果和伤亡须分开核对。对于医疗救助、慈善网络、普通人的伤亡与案狱处境，材料往往尤为稀薄。译写候选以编年材料定位；实录记载反映朝廷选择，崇祯期须另以奏疏、地方及敌对政权材料交叉。

\eventanchor{ML-1389-002}{明·1389（洪武二十二年）十二月}
\paragraph{1389（洪武二十二年）十二月：明蒙战争：司伦法投降明朝} 此一账本锚点将冲突、风险或法律处置置于具体时段。它所能可靠支持的判断是编年候选的年序可由官修与后出材料交叉；具体月日、参与者、战果或数字必须回查所列材料，不能将单一叙事当作定论。材料链由，因而同时保存了命令、报送、传记、地方记忆或跨政权叙事的不同立场。问题形成的条件是前期边防、地方武装或跨区域资源调动的压力推动了该行动；随后产生的可见影响是它改变了后续调兵、补给或地方秩序的条件；战果和伤亡须分开核对。对于医疗救助、慈善网络、普通人的伤亡与案狱处境，材料往往尤为稀薄。译写候选以编年材料定位；实录记载反映朝廷选择，崇祯期须另以奏疏、地方及敌对政权材料交叉。

\eventanchor{ML-1390-001}{明·1390（洪武二十三年）四月}
\paragraph{1390（洪武二十三年）四月：纳伊尔布哈和耀珠向明军投降} 此一账本锚点将冲突、风险或法律处置置于具体时段。它所能可靠支持的判断是编年候选的年序可由官修与后出材料交叉；具体月日、参与者、战果或数字必须回查所列材料，不能将单一叙事当作定论。材料链由，因而同时保存了命令、报送、传记、地方记忆或跨政权叙事的不同立场。问题形成的条件是前期边防、地方武装或跨区域资源调动的压力推动了该行动；随后产生的可见影响是它改变了后续调兵、补给或地方秩序的条件；战果和伤亡须分开核对。对于医疗救助、慈善网络、普通人的伤亡与案狱处境，材料往往尤为稀薄。译写候选以编年材料定位；实录记载反映朝廷选择，崇祯期须另以奏疏、地方及敌对政权材料交叉。

\eventanchor{ML-1391-001}{明·1391（洪武二十四年）五月}
\paragraph{1391（洪武二十四年）五月：阿贾施里叛乱并被镇压} 此一账本锚点将冲突、风险或法律处置置于具体时段。它所能可靠支持的判断是编年候选的年序可由官修与后出材料交叉；具体月日、参与者、战果或数字必须回查所列材料，不能将单一叙事当作定论。材料链由，因而同时保存了命令、报送、传记、地方记忆或跨政权叙事的不同立场。问题形成的条件是继承、官僚协调或皇权运作的压力使问题进入朝廷程序；随后产生的可见影响是它影响后续人事与政策议程；动机和派系标签不能由单一叙事裁定。对于医疗救助、慈善网络、普通人的伤亡与案狱处境，材料往往尤为稀薄。译写候选以编年材料定位；实录记载反映朝廷选择，崇祯期须另以奏疏、地方及敌对政权材料交叉。

\subsection{案狱、律令与地方程序}

\eventanchor{ML-1391-002}{明·1391（洪武二十四年）}
\paragraph{1391（洪武二十四年）：明军短暂占领哈密并撤退} 此一账本锚点将冲突、风险或法律处置置于具体时段。它所能可靠支持的判断是编年候选的年序可由官修与后出材料交叉；具体月日、参与者、战果或数字必须回查所列材料，不能将单一叙事当作定论。材料链由，因而同时保存了命令、报送、传记、地方记忆或跨政权叙事的不同立场。问题形成的条件是前期边防、地方武装或跨区域资源调动的压力推动了该行动；随后产生的可见影响是它改变了后续调兵、补给或地方秩序的条件；战果和伤亡须分开核对。对于医疗救助、慈善网络、普通人的伤亡与案狱处境，材料往往尤为稀薄。译写候选以编年材料定位；实录记载反映朝廷选择，崇祯期须另以奏疏、地方及敌对政权材料交叉。

\eventanchor{ML-1392-001}{明·1392（洪武二十五年）八月5日}
\paragraph{1392（洪武二十五年）八月5日：李松桂废黜王耀，成为朝鲜太祖；高句丽就此结束} 此一账本锚点将冲突、风险或法律处置置于具体时段。它所能可靠支持的判断是编年候选的年序可由官修与后出材料交叉；具体月日、参与者、战果或数字必须回查所列材料，不能将单一叙事当作定论。材料链由，因而同时保存了命令、报送、传记、地方记忆或跨政权叙事的不同立场。问题形成的条件是朝贡名义、边境安全与港口贸易的利益使交涉持续发生；随后产生的可见影响是它为后续册封、互市或海防安排留下文书与跨政权比较线索。对于医疗救助、慈善网络、普通人的伤亡与案狱处境，材料往往尤为稀薄。译写候选以编年材料定位；实录记载反映朝廷选择，崇祯期须另以奏疏、地方及敌对政权材料交叉。

\eventanchor{ML-1393-001}{明·1393（洪武二十六年）}
\paragraph{1393（洪武二十六年）：明军攻陷哈密} 此一账本锚点将冲突、风险或法律处置置于具体时段。它所能可靠支持的判断是编年候选的年序可由官修与后出材料交叉；具体月日、参与者、战果或数字必须回查所列材料，不能将单一叙事当作定论。材料链由，因而同时保存了命令、报送、传记、地方记忆或跨政权叙事的不同立场。问题形成的条件是前期边防、地方武装或跨区域资源调动的压力推动了该行动；随后产生的可见影响是它改变了后续调兵、补给或地方秩序的条件；战果和伤亡须分开核对。对于医疗救助、慈善网络、普通人的伤亡与案狱处境，材料往往尤为稀薄。译写候选以编年材料定位；实录记载反映朝廷选择，崇祯期须另以奏疏、地方及敌对政权材料交叉。

\eventanchor{ML-1394-001}{明·1394（洪武二十七年）}
\paragraph{1394（洪武二十七年）：明朝鲜朝贡关系正常化} 此一账本锚点将冲突、风险或法律处置置于具体时段。它所能可靠支持的判断是编年候选的年序可由官修与后出材料交叉；具体月日、参与者、战果或数字必须回查所列材料，不能将单一叙事当作定论。材料链由，因而同时保存了命令、报送、传记、地方记忆或跨政权叙事的不同立场。问题形成的条件是朝贡名义、边境安全与港口贸易的利益使交涉持续发生；随后产生的可见影响是它为后续册封、互市或海防安排留下文书与跨政权比较线索。对于医疗救助、慈善网络、普通人的伤亡与案狱处境，材料往往尤为稀薄。译写候选以编年材料定位；实录记载反映朝廷选择，崇祯期须另以奏疏、地方及敌对政权材料交叉。

\eventanchor{ML-1394-002}{明·1394（洪武二十七年）}
\paragraph{1394（洪武二十七年）：户部、布政司与里甲体系围绕户籍、田土和赋役的申报形成年度行政链，本条标记其可回查节点} 此一账本锚点将冲突、风险或法律处置置于具体时段。它所能可靠支持的判断是来源导向的年度桥接条目：本年材料可定位相关政务线索；具体月日、发文机关、地域和执行结果仍须逐项回查，不能以制度文本或奏报替代实际效果。材料链由，因而同时保存了命令、报送、传记、地方记忆或跨政权叙事的不同立场。问题形成的条件是财政征解、交通工程或货币支付的压力使相关措施进入政务；随后产生的可见影响是其法定安排不等于地方执行，后续须以地域材料追踪负担与成效。对于医疗救助、慈善网络、普通人的伤亡与案狱处境，材料往往尤为稀薄。桥接条目只标示可回查的年度问题与材料链；实录有中央选择性，崇祯期尤其需要奏疏、地方、朝鲜及满文材料交叉。

\eventanchor{ML-1395-001}{明·1395（洪武二十八年）}
\paragraph{1395（洪武二十八年）：户部、布政司与里甲体系围绕户籍、田土和赋役的申报形成年度行政链，本条标记其可回查节点} 此一账本锚点将冲突、风险或法律处置置于具体时段。它所能可靠支持的判断是来源导向的年度桥接条目：本年材料可定位相关政务线索；具体月日、发文机关、地域和执行结果仍须逐项回查，不能以制度文本或奏报替代实际效果。材料链由，因而同时保存了命令、报送、传记、地方记忆或跨政权叙事的不同立场。问题形成的条件是财政征解、交通工程或货币支付的压力使相关措施进入政务；随后产生的可见影响是其法定安排不等于地方执行，后续须以地域材料追踪负担与成效。对于医疗救助、慈善网络、普通人的伤亡与案狱处境，材料往往尤为稀薄。桥接条目只标示可回查的年度问题与材料链；实录有中央选择性，崇祯期尤其需要奏疏、地方、朝鲜及满文材料交叉。

\subsection{边地冲突与行政后果}

\eventanchor{ML-1395-002}{明·1395（洪武二十八年）}
\paragraph{1395（洪武二十八年）：北元诸部、东北和西南的边报、招抚与防御命令持续进入本年政务，须按具体部众和地域复核} 此一账本锚点将冲突、风险或法律处置置于具体时段。它所能可靠支持的判断是来源导向的年度桥接条目：本年材料可定位相关政务线索；具体月日、发文机关、地域和执行结果仍须逐项回查，不能以制度文本或奏报替代实际效果。材料链由，因而同时保存了命令、报送、传记、地方记忆或跨政权叙事的不同立场。问题形成的条件是朝贡名义、边境安全与港口贸易的利益使交涉持续发生；随后产生的可见影响是它为后续册封、互市或海防安排留下文书与跨政权比较线索。对于医疗救助、慈善网络、普通人的伤亡与案狱处境，材料往往尤为稀薄。桥接条目只标示可回查的年度问题与材料链；实录有中央选择性，崇祯期尤其需要奏疏、地方、朝鲜及满文材料交叉。

\eventanchor{ML-1396-001}{明·1396（洪武二十九年）四月}
\paragraph{1396（洪武二十九年）四月：明军击败勃林帖木儿} 此一账本锚点将冲突、风险或法律处置置于具体时段。它所能可靠支持的判断是编年候选的年序可由官修与后出材料交叉；具体月日、参与者、战果或数字必须回查所列材料，不能将单一叙事当作定论。材料链由，因而同时保存了命令、报送、传记、地方记忆或跨政权叙事的不同立场。问题形成的条件是前期边防、地方武装或跨区域资源调动的压力推动了该行动；随后产生的可见影响是它改变了后续调兵、补给或地方秩序的条件；战果和伤亡须分开核对。对于医疗救助、慈善网络、普通人的伤亡与案狱处境，材料往往尤为稀薄。译写候选以编年材料定位；实录记载反映朝廷选择，崇祯期须另以奏疏、地方及敌对政权材料交叉。

\eventanchor{ML-1396-002}{明·1396（洪武二十九年）}
\paragraph{1396（洪武二十九年）：学校、科举、礼制或律令的颁行与讨论在本年材料中留下制度线索，规定与地方实践应分开处理} 此一账本锚点将冲突、风险或法律处置置于具体时段。它所能可靠支持的判断是来源导向的年度桥接条目：本年材料可定位相关政务线索；具体月日、发文机关、地域和执行结果仍须逐项回查，不能以制度文本或奏报替代实际效果。材料链由，因而同时保存了命令、报送、传记、地方记忆或跨政权叙事的不同立场。问题形成的条件是制度、教育或知识传播的需要推动了相关行动或文本形成；随后产生的可见影响是它提供制度和传播史的锚点，不能据此推断社会整体接受。对于医疗救助、慈善网络、普通人的伤亡与案狱处境，材料往往尤为稀薄。桥接条目只标示可回查的年度问题与材料链；实录有中央选择性，崇祯期尤其需要奏疏、地方、朝鲜及满文材料交叉。

\eventanchor{ML-1397-001}{明·1397（洪武三十年）十月}
\paragraph{1397（洪武三十年）十月：林宽叛乱：锦族叛乱被击败} 此一账本锚点将冲突、风险或法律处置置于具体时段。它所能可靠支持的判断是编年候选的年序可由官修与后出材料交叉；具体月日、参与者、战果或数字必须回查所列材料，不能将单一叙事当作定论。材料链由，因而同时保存了命令、报送、传记、地方记忆或跨政权叙事的不同立场。问题形成的条件是前期边防、地方武装或跨区域资源调动的压力推动了该行动；随后产生的可见影响是它改变了后续调兵、补给或地方秩序的条件；战果和伤亡须分开核对。对于医疗救助、慈善网络、普通人的伤亡与案狱处境，材料往往尤为稀薄。译写候选以编年材料定位；实录记载反映朝廷选择，崇祯期须另以奏疏、地方及敌对政权材料交叉。

\eventanchor{ML-1397-002}{明·1397（洪武三十年）十二月}
\paragraph{1397（洪武三十年）十二月：明蒙毛干预：斯伦法被废黜并请求明援助以恢复其权力} 此一账本锚点将冲突、风险或法律处置置于具体时段。它所能可靠支持的判断是编年候选的年序可由官修与后出材料交叉；具体月日、参与者、战果或数字必须回查所列材料，不能将单一叙事当作定论。材料链由，因而同时保存了命令、报送、传记、地方记忆或跨政权叙事的不同立场。问题形成的条件是继承、官僚协调或皇权运作的压力使问题进入朝廷程序；随后产生的可见影响是它影响后续人事与政策议程；动机和派系标签不能由单一叙事裁定。对于医疗救助、慈善网络、普通人的伤亡与案狱处境，材料往往尤为稀薄。译写候选以编年材料定位；实录记载反映朝廷选择，崇祯期须另以奏疏、地方及敌对政权材料交叉。

\eventanchor{ML-1398-001}{明·1398（洪武三十一年）一月}
\paragraph{1398（洪武三十一年）一月：明蒙毛干预：斯伦法恢复权力} 此一账本锚点将冲突、风险或法律处置置于具体时段。它所能可靠支持的判断是编年候选的年序可由官修与后出材料交叉；具体月日、参与者、战果或数字必须回查所列材料，不能将单一叙事当作定论。材料链由，因而同时保存了命令、报送、传记、地方记忆或跨政权叙事的不同立场。问题形成的条件是继承、官僚协调或皇权运作的压力使问题进入朝廷程序；随后产生的可见影响是它影响后续人事与政策议程；动机和派系标签不能由单一叙事裁定。对于医疗救助、慈善网络、普通人的伤亡与案狱处境，材料往往尤为稀薄。译写候选以编年材料定位；实录记载反映朝廷选择，崇祯期须另以奏疏、地方及敌对政权材料交叉。
